\section{Evaluation, datasets and benchmarks}
\label{sec:evaluation}

Probes, concept directions, dictionaries, attention patterns and steered outputs have no self-evident
correctness. This section covers what is evaluated, the data and concept resources available, the
benchmarks and protocols in use, and what limits comparison between studies. Every count here and in \S\ref{sec:findings} refers to the \Ncore{} core studies and \Nclaims{}
records of \S\ref{sec:method}.

\subsection{What is evaluated}
\label{sec:eval-what}

Five non-interchangeable properties recur.

\textbf{Faithfulness and probe validity.} A faithful artefact reflects the computation the model
performs~\cite{jacovi2020towards}, which for probes, directions and features means connecting it to the
output through intervention (\S\ref{sec:intervention}); how a removal procedure is chosen already reorders
attribution methods~\cite{hooker2019benchmark}, and fidelity and sensitivity measures are themselves
choices~\cite{yeh2019fidelity}. Probe validity is the weaker property that performance
reflects the representation rather than the decoder's own work, tested by control-task selectivity,
description-length probes or shuffled labels~\cite{hewitt2019designing,voita2020information}, with a
randomly initialised encoder calibrating what the architecture alone affords. Perturbation and model
randomisation test attribution and attention~\cite{adebayo2018sanity}, which medical saliency and attention
maps can fail~\cite{arun2021trustworthiness,sadanandan2026attention}.

\textbf{Concept validity.} Concepts from label schemas, reports or language models may be
invalid; evaluation covers agreement with expert annotation, inter-rater reliability, whether the concept
set is complete enough to account for the prediction~\cite{yeh2019completeness}, concept-bank coverage of
clinical vocabulary and, for post-hoc names, agreement between descriptions and reports judged by a frozen
medical language model~\cite{haque2026medconcept}. Plausibility, whether a clinician finds an
explanation reasonable, is a property of the reader and independent of
faithfulness~\cite{jacovi2020towards,miller2017explanation}.

\textbf{Stability.} Discovered features and learned directions depend on seeds, dictionary size,
training subsets and model versions; reproducibility across seeds, agreement of directions across
checkpoints and consistency of probing across splits measure it, and the studies reporting them find
substantial instability (\S\ref{sec:discovery}).

\textbf{Specificity of interventions.} An intervention is evaluated against the three causal claims
of \S\ref{sec:methods}: causal sensitivity, direction-selective sensitivity against matched controls,
and concept-specific causal contribution (\S\ref{sec:intervention}).

\textbf{Clinical utility.} The properties above are functionally grounded; utility is the
application-grounded end of the same scale~\cite{doshivelez2017towards}. An artefact has clinical utility
when it helps a clinician or developer decide better, measured by reader studies, task-level end-points such as shortcut detection or error
correction, and downstream metrics after intervention. The corpus has one dermatology reader study
(\S\ref{sec:eval-bench}), and systematic reviews find mixed, sometimes bidirectional effects on trust and
performance~\cite{rosenbacke2024how}.

\textbf{How often the five properties are reported.} Unevenly, and Table~\ref{tab:vcov} gives the
count for every class of check. Three gaps stand out. Probe validity is almost never tested: a
permuted-label null or chance calibration appears in \Nprobectrl{} of the \Ndecni{} non-intrinsic decoding
studies and a randomly initialised encoder in \Nproberand{}, and no study runs a randomized control task.
Stability is reported by \Nsaestab{} of the \Nsae{} dictionary studies and cross-seed identity of features
or names by none, nor does any decoding, geometry or intervention study establish it for a probe or
direction. And no core study assesses clinical utility with a reader study, the one dermatology study
being contextual, although clinician-facing explanation has been argued to fail exactly at this step,
where an artefact that looks reasonable is taken for one that is faithful~\cite{ghassemi2021falsehope}.
The corpus therefore rests on uncontrolled decodability and unvalidated or singly validated feature names
(\S\ref{sec:findings}).

\subsection{Datasets and concept resources}
\label{sec:eval-data}

%% [written after the `evaluation' bucket is in; Table 10]
\input{sections/sec5_2_datasets}

\subsection{Benchmarks and protocols}
\label{sec:eval-bench}

%% [written after the `evaluation' bucket is in]
\textbf{Benchmarks that evaluate representations.} A few benchmarks score frozen medical features rather than fine-tuned
performance. Pathology benchmarks compare 19 models on
13 external cohorts~\cite{neidlinger2025benchmark} and 23 models on 16 tile-level datasets with their feature
spaces~\cite{marza2025thunder}; others test clinical end-points~\cite{campanella2024clinical},
standardise slide processing~\cite{zhang2025pathobench}, probe spatial
transcriptomics~\cite{jaume2024hest} or weigh biology against centre
confounders~\cite{dejong2025current}. In radiology, BenchX covers nine
chest-radiograph datasets~\cite{zhou2024benchx}; frozen-feature comparisons cover medical and general
chest-radiograph encoders~\cite{li2025featurequality}, ten three-dimensional CT encoders on three
thoracic cohorts~\cite{chevli2026frozenct} and general self-supervised features in more than 200
radiology settings~\cite{baharoon2023dinov2}. Representation benchmarks require controls: an EEG
benchmark uses label permutation, scrambled-label fine-tuning and random-init
encoders~\cite{zare2026eeg}, an ECG benchmark matches dictionary capacity across depths, widths and
seeds~\cite{duan2026ecginterpbench}; both lie outside medical imaging, where such standards are mostly
absent.

\textbf{Faithfulness protocols.} Model- and data-randomisation checks~\cite{adebayo2018sanity} became a
four-criterion radiology protocol that eight saliency methods failed~\cite{arun2021trustworthiness}. A
radiologist-segmentation benchmark~\cite{saporta2022benchmarking} and revised
criteria~\cite{zhang2023revisiting} followed. Medical vision--language audits find attention unfaithful
under occlusion~\cite{sadanandan2026attention}, low grounding accuracy~\cite{chen2026grounding} and weak
reasoning faithfulness under perturbation~\cite{moll2025faithfulness}. A test-bed builds faithful and unfaithful pairs by model editing to test the metrics
themselves~\cite{zaman2025causal}.

\textbf{Feature and intervention evaluation.} Sparse-autoencoder evaluation includes SAEBench and its re-seeding audit, SynthSAEBench's synthetic ground truth, cross-seed feature sharing and intervention-scored interpretation~\cite{karvonen2025saebench,chanin2026are,chanin2026synthsaebench,paulo2025sparse,paulo2024automatically} (\S\ref{sec:discovery}). A unified steering framework covers multiple-choice and open-ended
settings~\cite{im2025unifiedsteering}; reliability, non-identifiability and safety studies provide
negative controls (\S\ref{sec:intervention})~\cite{tan2024steeringreliability,venkatesh2026nonident,korznikov2025the}.

\textbf{Tooling.} The designs of \S\ref{sec:methods} run on open software: sparse dictionaries with the
SAELens and dictionary-learning libraries, patching and
steering of open multimodal LLMs with TransformerLens, nnsight or pyvene, and closed-form concept erasure
with a reference LEACE implementation~\cite{belrose2023leace}. Released medical
dictionaries~\cite{abdulaal2024an,nooralahzadeh2026universal} and the checkpoints of
Table~\ref{tab:models} supply the objects. The four encoders with embeddings but not weights
(\S\ref{sec:models}) expose no intermediate activations, admitting decoding and geometry but not tracing
or within-run intervention.

\textbf{Human evaluation and reporting standards.} A three-phase study found dermatologist-aligned
explanations increased trust and confidence~\cite{chanda2023dermatologist}; systematic reviews find few eligible studies, mixed effects
and high risk of bias~\cite{rosenbacke2024how,uddin2026which,gambetti2025a}, supporting no claim of
clinical benefit. Guidelines cover clinical XAI
evaluation~\cite{jin2022guidelines,chen2021humancentered}, explanation
properties~\cite{nauta2023interpreting}, user-centred studies~\cite{rong2022userstudies,donosoguzman2025a}
and mechanistic-interpretability audits~\cite{lan2026auditable}, none addressing representation-level
claims about medical foundation models.


\subsection{Protocol comparability and evidence quality}
\label{sec:eval-compare}

%% [written after the corpus is coded; Table 11]
Published values are rarely commensurable, so side-by-side display is inappropriate;
Table~\ref{tab:famreq} states what each family requires and how far this corpus has taken it; the counts of comparable evaluations are below. Decoding protocols
overlap most, and only within a modality and a shared public label schema; probes otherwise report AUROC
on different concept sets, splits, label definitions, decoder capacities and controls, obscuring
differences of a few points. Discovery overlaps least: width, sparsity, layer and
interpretation vary, cross-seed stability is seldom reported and nominal concepts cannot be matched.
Intervention records permit individual inspection, but differing end-points, magnitudes and controls preclude
effect-size comparison. The most compatible stratum, six chest-radiograph studies using linear decoders on CheXpert or MIMIC-CXR, still differs in targets, label versions, splits and
regularisation, so it shows six distinct estimands rather than a poolable stratum, which the reporting
checklists of \S\ref{sec:future} would supply. Independent replication by a second group is uncoded and
apparently rare, and same-data medical--general comparisons confound domain with data, objective,
architecture and protocol (\S\ref{sec:findings}).

\textbf{Modality and family.} Pathology appears in \Npathstud{} of the \Ncore{} studies and chest radiography in
\Ncxrstud{} (\Npathrec{} and \Ncxrrec{} records); every other modality is markedly rarer
(Fig.~\ref{fig:corpus}c). Chest radiography supplies \Nintcxr{} of the \Nintrec{} intervention records and
\Ntracecxr{} of the \Ntracerec{} tracing records, so multimodal-LLM interpretability is largely a
chest-radiograph literature; pathology dominates geometry and discovery, entering intervention only through
frozen-embedding post-processing. Modality is retained in the evidence tables because a claim rarely
transfers across modalities.

\textbf{Study quality.} Study-quality coding comes from full texts (four from abstracts,
\S\ref{sec:method}). The fields are the subset of the imaging-AI and clinical-AI reporting
checklists~\cite{mongan2020claim,norgeot2020minimum} that a representation-level claim can be scored on from a paper: they
record what was reported, not whether it was done well.
Of the \Ncore{} core studies, \Npreprint{} are preprints, \Nconf{} conference
papers, \Nwork{} workshop papers and \Njour{} journal articles; \Ngrey{} of the preprints are grey-literature deposits and
\Nsingle{} are single-author, admitted under the rule of \S\ref{sec:method} that any retrievable full text meeting the eligibility criteria is coded. Excluding the
grey deposits lowers the corpus-breadth tier of F8a$''$ alone and changes no direction; excluding all \Ngreysingle{} lowers F3c as well. Of these studies, \Nextyes{} evaluate on a cohort or
dataset unused for probe, dictionary or intervention training, \Nextno{} do not and \Nextunk{} are
indeterminate;
patient-, slide- or subject-level splits are stated in \Nsplityes{} and unstated in \Nsplitno{}. Seeds,
repeated runs, intervals or tests accompany the main result in \Nseedyes{}, not in \Nseedno{}, and
indeterminately in \Nseedunk{}. Of the \Nclaims{} records, \Ndirpos{} report a positive result,
\Ndirmix{} a mixed or partial one and \Ndirnull{} a null or negative one; null results this rare warrant
publication- and reporting-bias caution for every statement below. The young corpus relies more on internal
than external validation, a second reason \S\ref{sec:findings} grades breadth rather than pooling effect
sizes.

\textbf{Concept provenance and concept validity.} Among those records, \Nprovlabel{} use structured
schemas, \Nprovtext{} text-mined concepts, \Nprovexpert{} expert vocabularies, \Nprovmodel{}
model-generated banks and \Nprovnone{} no prior concept, \Nposthocname{} of the last naming features post
hoc. Expert-defined and model-generated concepts occur almost only in intrinsic designs
(\S\ref{sec:intrinsic}). The best-supported decoding findings (\S\ref{sec:findings}, F1a and F2) rest on
labels automatically extracted from chest-radiograph
reports~\cite{irvin2019chexpert,johnson2019mimicjpg}, so concept validity reflects a labeller, not clinician
judgement (\S\ref{sec:eval-data}). Discovery without a prior concept shifts the burden to post-hoc naming,
clinician-validated by fewer than half of dictionary studies (\S\ref{sec:eval-what}). Provenance locates where a claim may fail for reasons unrelated to the model.


